\chapter{Brownian Motion}
\section{Lecture 8 Continued}

Work on a filtered probability space
$(\Omega,\mathcal{F},(\mathcal{F}_t)_{t\ge0},\mathbb{P})$.

\begin{definition}[Standard Brownian motion]
An adapted process $B=(B_t)_{t\ge0}$ is a standard one-dimensional
$(\mathcal{F}_t)$-Brownian motion if:
\begin{enumerate}
 \item $B_0=0$ almost surely;
 \item $B$ has continuous sample paths almost surely;
 \item for $0\le s\le t$, the increment $B_t-B_s$ is independent of
 $\mathcal{F}_s$ and
 \[
 B_t-B_s\sim\mathcal{N}(0,t-s).
 \]
\end{enumerate}
\end{definition}

\begin{definition}[Gaussian process]
A process $(G_t)_{t\ge0}$ is called a \emph{Gaussian process} if, for every
finite collection of times $t_1,\ldots,t_n$, the random vector
$(G_{t_1},\ldots,G_{t_n})$ has a multivariate normal distribution.
\end{definition}

\begin{proposition}[Gaussian characterization of Brownian motion]
Let $B$ be an adapted process with continuous paths almost surely. Relative
to its natural filtration
\[
 \mathcal{F}_t^B:=\sigma(B_s:0\le s\le t),
\]
$B$ is Brownian motion if and only if:
\begin{enumerate}
 \item $B$ is a Gaussian process;
 \item $\mathbb{E}[B_t]=0$ for every $t\ge0$;
 \item $\mathbb{E}[B_sB_t]=s\wedge t$ for every $s,t\ge0$.
\end{enumerate}
\end{proposition}
\begin{proof}
Suppose first that $B$ is Brownian motion. Every finite vector
$(B_{t_1},\ldots,B_{t_n})$ is Gaussian because it can be written as a linear
transformation of the independent Gaussian increments. Also,
$\mathbb{E}[B_t]=0$. If $s\le t$, then
\begin{align*}
 \mathbb{E}[B_sB_t]
 &=\mathbb{E}\bigl[B_s(B_s+B_t-B_s)\bigr]\\
 &=\mathbb{E}[B_s^2]+\mathbb{E}[B_s(B_t-B_s)]\\
 &=s,
\end{align*}
since $B_t-B_s$ is independent of $\mathcal{F}_s$ and has mean zero.

Conversely, assume the three displayed properties. Since
$\mathbb{E}[B_0^2]=0$, we have $B_0=0$ almost surely. For $s\le t$,
the Gaussian variable $B_t-B_s$ has mean zero and variance
\[
 \operatorname{Var}(B_t-B_s)=t+s-2s=t-s.
\]
Hence $B_t-B_s\sim\mathcal{N}(0,t-s)$. Moreover, for every
$r\le s$,
\[
 \operatorname{Cov}(B_t-B_s,B_r)=(r\wedge t)-(r\wedge s)=r-r=0.
\]
The vector $(B_t-B_s,B_{r_1},\ldots,B_{r_m})$ is jointly Gaussian whenever
$r_1,\ldots,r_m\le s$. For jointly Gaussian random variables, zero
covariance implies independence. Thus $B_t-B_s$ is independent of
$\mathcal{F}_s^B$, proving the Brownian-increment property.
\end{proof}

\noindent\textit{Filtration caveat.} The preceding characterization proves
independence from the \emph{natural} filtration. If a larger filtration is
used, one must additionally assume that it contains no information about
future Brownian increments; otherwise the covariance conditions alone do not
imply independence from $\mathcal{F}_s$.

\begin{theorem}[L\'evy's characterization]
Let $X$ be an $(\mathcal{F}_t)$-adapted continuous process with $X_0=0$.
Then $X$ is an $(\mathcal{F}_t)$-Brownian motion if and only if both $X$ and
$(X_t^2-t)_{t\ge0}$ are martingales.
\end{theorem}

\begin{proof}[Proof of the ``only if'' direction]
If $X$ is Brownian motion, then it is a martingale because its increments
are independent of the past and have mean zero. For $s\le t$, write
$X_t=X_s+(X_t-X_s)$. Conditional on $\mathcal{F}_s$, the increment
$X_t-X_s$ has mean zero and second moment $t-s$. Therefore
\begin{align*}
 \mathbb{E}[X_t^2-t\mid\mathcal{F}_s]
 &=X_s^2+2X_s\mathbb{E}[X_t-X_s\mid\mathcal{F}_s]
   +\mathbb{E}[(X_t-X_s)^2\mid\mathcal{F}_s]-t\\
 &=X_s^2+(t-s)-t=X_s^2-s.
\end{align*}
Thus $(X_t^2-t)_{t\ge0}$ is a martingale.
\end{proof}

\begin{theorem}[Strong Markov property of Brownian motion]
Let $(B_t)_{t\ge0}$ be a Brownian motion with respect to a filtration
$(\mathcal{F}_t)_{t\ge0}$ satisfying the usual conditions, and let $\tau$ be
an almost surely finite stopping time. Then the shifted process
\[
W_t:=B_{\tau+t}-B_\tau,\qquad t\ge0,
\]
is a Brownian motion independent of $\mathcal{F}_\tau$.

Equivalently, if $0\le t_1<\cdots<t_n$ and $h:\mathbb{R}^n\to\mathbb{R}$
is bounded and Borel measurable, then
\[
\mathbb{E}\!\left[
h\bigl(B_{\tau+t_1}-B_\tau,\ldots,B_{\tau+t_n}-B_\tau\bigr)
\;\middle|\; \mathcal{F}_\tau
\right]
=
\mathbb{E}\!\left[h(B_{t_1},\ldots,B_{t_n})\right].
\]
\end{theorem}

\begin{proof}
We first prove the result when $\tau$ is a simple stopping time, taking
values $r_1,\ldots,r_m$. Let $A\in\mathcal{F}_\tau$. Since
$A\cap\{\tau=r_j\}\in\mathcal{F}_{r_j}$, the independent-increments
property of Brownian motion gives
\[
\begin{aligned}
&\mathbb{E}\!\left[
\mathbf{1}_A
h\bigl(B_{\tau+t_1}-B_\tau,\ldots,B_{\tau+t_n}-B_\tau\bigr)
\right] \\
&\qquad =
\sum_{j=1}^m
\mathbb{E}\!\left[
\mathbf{1}_{A\cap\{\tau=r_j\}}
h\bigl(B_{r_j+t_1}-B_{r_j},\ldots,B_{r_j+t_n}-B_{r_j}\bigr)
\right] \\
&\qquad =
\sum_{j=1}^m
\mathbb{P}(A\cap\{\tau=r_j\})
\mathbb{E}\!\left[h(B_{t_1},\ldots,B_{t_n})\right] \\
&\qquad =
\mathbb{P}(A)\mathbb{E}\!\left[h(B_{t_1},\ldots,B_{t_n})\right].
\end{aligned}
\]
This is exactly the desired conditional-expectation identity.

For a general bounded stopping time $\tau$, approximate it from above by
simple stopping times:
\[
\tau_n:=2^{-n}\left\lceil 2^n\tau\right\rceil.
\]
Then $\tau_n\downarrow\tau$. Since Brownian paths are continuous,
\[
B_{\tau_n+t_i}-B_{\tau_n}
\longrightarrow
B_{\tau+t_i}-B_\tau
\qquad\text{almost surely}
\]
for each $i$. Applying the simple-stopping-time result to $\tau_n$ and
passing to the limit using bounded convergence proves the result for bounded
$\tau$. An arbitrary almost surely finite stopping time follows by applying
the bounded case to $\tau\wedge N$ and then letting $N\to\infty$.
\end{proof}

\begin{remark}[Interpretation]
The ordinary Markov property says that, after a fixed deterministic time
$s$, the future process $(B_{s+t}-B_s)_{t\ge0}$ is a new Brownian motion
independent of the past $\mathcal{F}_s$.

The strong Markov property says the same remains true if the time $s$ is
replaced by a random time $\tau$, provided that $\tau$ is a stopping time:
the decision that ``now is time $\tau$'' may use information from the past,
but not information from the future.
\end{remark}

\begin{theorem}[Reflection principle]
Let $a>0$ and define
\[
T_a:=\inf\{s\ge0:B_s=a\},
\qquad
M_t:=\sup_{0\le s\le t}B_s.
\]
Then, for every $t>0$,
\[
\mathbb{P}(M_t\ge a)
=
\mathbb{P}(T_a\le t)
=
2\mathbb{P}(B_t\ge a).
\]
Consequently,
\[
M_t\overset{d}{=}|B_t|.
\]
\end{theorem}

\begin{proof}
Since Brownian paths are continuous and $B_0=0<a$,
\[
\{M_t\ge a\}=\{T_a\le t\}.
\]

We now explain the reflection argument. Consider a path for which
$T_a\le t$ and $B_t<a$. Such a path reaches the level $a$ for the first
time at $T_a$, and then ends below $a$. Reflect the part of the path after
this first hitting time across the horizontal line at height $a$:
\[
\widetilde B_s:=
\begin{cases}
B_s, & 0\le s\le T_a,\\
2a-B_s, & s>T_a.
\end{cases}
\]
The reflected path has the same history up to $T_a$, but its terminal value is
\[
\widetilde B_t=2a-B_t>a.
\]

Why does this reflection preserve probabilities? By the strong Markov
property, conditional on $\mathcal{F}_{T_a}$, the post-hitting process
\[
W_u:=B_{T_a+u}-B_{T_a},\qquad u\ge0,
\]
is a fresh Brownian motion independent of $\mathcal{F}_{T_a}$. Brownian
motion is symmetric, meaning
\[
(W_u)_{u\ge0}\overset{d}{=}(-W_u)_{u\ge0}.
\]
But reflecting the original path after $T_a$ replaces $W$ by $-W$, because
\[
\widetilde B_{T_a+u}-a
=
2a-B_{T_a+u}-a
=
-(B_{T_a+u}-a)
=
-W_u.
\]
Hence reflection preserves the conditional distribution of the future given
the past up to $T_a$.

Reflecting twice gives back the original path. Thus reflection pairs paths in
\[
\{T_a\le t,\ B_t<a\}
\]
with paths in
\[
\{B_t>a\}.
\]
The latter paths necessarily hit $a$ before time $t$, since they start at
$0<a$ and end above $a$. Therefore,
\[
\mathbb{P}(T_a\le t,\ B_t<a)=\mathbb{P}(B_t>a).
\]

Now split according to the terminal value:
\[
\begin{aligned}
\mathbb{P}(T_a\le t)
&=
\mathbb{P}(T_a\le t,\ B_t\ge a)
+
\mathbb{P}(T_a\le t,\ B_t<a) \\
&=
\mathbb{P}(B_t\ge a)+\mathbb{P}(B_t>a).
\end{aligned}
\]
Since $B_t\sim\mathcal{N}(0,t)$ has a continuous distribution,
\[
\mathbb{P}(B_t>a)=\mathbb{P}(B_t\ge a).
\]
Hence
\[
\mathbb{P}(T_a\le t)=2\mathbb{P}(B_t\ge a).
\]

Finally, for $x\ge0$,
\[
\begin{aligned}
\mathbb{P}(M_t\le x)
&=1-\mathbb{P}(M_t>x) \\
&=1-2\mathbb{P}(B_t>x) \\
&=\mathbb{P}(|B_t|\le x).
\end{aligned}
\]
Therefore $M_t\overset{d}{=}|B_t|$.
\end{proof}
\section{Lecture 9}

Some scaling properties of Brownian motion (self-similarity):
\[
\{B_{at}\}_{t \ge 0} \overset{d}{=} \{a^{1/2} B_t\}_{t \ge 0}
\]
Insert proof (just algebra)

\begin{theorem}
Let $B$ be a BM. Then, the process
\[
\begin{cases}
X_0 &= 0 \\
X_t &= tB_{1/t} \quad t > 0
\end{cases}
\]
is also a BM.
\end{theorem}
For the proof, we check the algebra for the second definition of BM.

\begin{corollary}
$\Prob( \inf\{
    t > 0 \colon B_t = 0
\} = 0) = 1$. Ie, 0 is almost surely an accumulation point of zeroes of a BM. In fact, any zero is an accumulation point.
\end{corollary}

\begin{theorem}
Almost surely, the path of a BM is of unbounded variation on any compact interval. (This is why Stochastic Calculus exists!!)
\end{theorem}
\begin{proof}
$\{\Pi_n\}$ be an increasing sequence of partitions of $[a,b]$. Further, let $\{\Pi_n\}$ be such that the mesh of $\{\Pi_n\}$ goes to zero. Now, consider
\begin{align*}
&\left(\sum_{t_i \in \Pi_n} (B_{t_{i+1}}-B_{t_i})^2\right) - (b-a) \\
&= \sum_{t_i \in \Pi_n} (B_{t_{i+1}}-B_{t_i})^2 - (t_{i+1}-t_i) \\
\end{align*}
Now, we will show that this goes to zero in $L^2$. All the terms within the sum are independent random variables with mean zero. This then shows that
\begin{align*}
&\E\left((\sum_{t_i \in \Pi_n} (B_{t_{i+1}}-B_{t_i})^2 - (t_{i+1}-t_i))^2\right) \\
&= \sum_{t_i \in \Pi_n} \E((B_{t_{i+1}}-B_{t_i})^2 - (t_{i+1}-t_i))^2
\end{align*}
We know that
\[
\frac{(B_{t_{i+1}}-B_{t_i})^2}{t_{i+1}-t_i} \sim Z^2,
\]
where $Z \sim \mcal{N}(0,1)$. So, we then get
\begin{align*}
&\sum_{t_i \in \Pi_n} (t_{i+1}-t_i)^2 \E((Z^2-1)^2) \\
& \le \E((Z^2-1)^2)  \sup_{t_i \in \Pi_n} |t_{i+1}-t_i| (b-a) \\
&\to 0,
\end{align*}
since the mesh of $\{\Pi_n\}$ goes to zero.

So, $\sum_{t_i \in \Pi_n} (B_{t_{i+1}} - B_{t_i})^2 \overset{L^2}{\to} b-a$.

Thus, there exists a subsequence $\{\Pi_n'\}$ such that we have almost sure convergence. On the other hand,
\begin{align*}
&\sum_{t_i \in \Pi_n'} (B_{t_{i+1}}-B_{t_i})^2 \\
&\le \sup_{t_i \in \Pi_n'} |B_{t_{i+1}}-B_{t_i}| \cdot \sum_{t_i \in \Pi_n'}|B_{t_{i+1}}-B_{t_i}| \\
&\le \sup_{t_i \in \Pi_n'} |B_{t_{i+1}}-B_{t_i}|  \cdot \operatorname{TV}_{[a,b]}(B)
\end{align*}
Now, since $B$ is continuous on a compact interval, we can use uniform continuity, so that $\sup_{t_i \in \Pi_n'} |B_{t_{i+1}}-B_{t_i}|  \to 0$. This means that the total variation cannot be finite, since if it was, the summation expression would converge to 0, which would not have $L^2$ norm going to $b-a$.
\end{proof}
